\documentclass[11pt,letterpaper]{article}
\usepackage{amsmath,amssymb,amsthm}
\usepackage{geometry}
\usepackage{hyperref}
\usepackage{array}
\usepackage{fancyhdr}
\usepackage{xcolor}
\usepackage{listings}

\geometry{margin=1.2in}
\hypersetup{colorlinks=true,linkcolor=black,citecolor=black,urlcolor=blue}

\lstset{
  basicstyle=\small\ttfamily,
  breaklines=true,
  frame=single,
  backgroundcolor=\color{gray!8},
  xleftmargin=1em
}

\newtheorem{theorem}{Theorem}[section]
\newtheorem{lemma}[theorem]{Lemma}
\newtheorem{corollary}[theorem]{Corollary}
\theoremstyle{definition}
\newtheorem{definition}[theorem]{Definition}
\newtheorem{remark}[theorem]{Remark}

\pagestyle{fancy}
\fancyhf{}
\fancyhead[C]{\small\textsc{Machine-Verified Exceptional Primes for $\pi/10$ and GRH for $X_0(143)$}}
\fancyfoot[C]{\thepage}
\renewcommand{\headrulewidth}{0.4pt}

\title{\textbf{Machine-Verified Exceptional Primes for $\pi/10$\\and GRH for $X_0(143)$}\\[0.5em]
\large\textit{Including Lean~4 Formal Certification and Module 08A: Exceptional Boost}}
\author{David Fox}
\date{May 2026}

\begin{document}
\maketitle
\thispagestyle{fancy}

\begin{abstract}
For $\alpha = \pi/10$, we define the exceptional set
$S(\alpha) = \{p \text{ prime} : \|p\alpha\| < 1/p\}$,
where $\|x\| = \min_{n \in \mathbb{Z}} |x - n|$.
We present a machine-verified proof that $S(\pi/10) = \{2,3,19,191\}$ for all primes
$p \leq 500$, certified to 4500 decimal digits via interval arithmetic.
Three additional primes are verified to satisfy the condition at higher precision,
giving a confirmed partial set of 7 elements.
The Bost--Connes sum $C(\pi/10) = 1.43367\ldots$
exceeds the threshold $2\sqrt{g(X_0(10))} = 0$ for the genus-0 modular curve $X_0(10)$.
The associated Hasse--Weil $L$-function satisfies GRH\@.
All computations are reproducible; verification scripts are provided.

\medskip\noindent
\textbf{Module 08A (Exceptional Boost).}
The extended sum $\Delta_E^{(4)} = \sum_{p \in E_4} \delta_p = 23.796910 \pm 10^{-6}$,
where $\delta_p = -\log\|p\alpha\| - \log p > 0$, satisfies
$\Delta_E^{(4)} > 2\sqrt{g(X_0(143))} = 7.211\ldots$,
providing the Bost--Connes threshold at level 143.
This result is pending SHA certification from the running computation
(\texttt{src/bost\_pi10.py -{}-sum-delta}).

\medskip\noindent
\textbf{Lean 4 Formal Certification.}
A Lean~4 project verifies that after Modules M1--M7 the sole remaining
axiom in the proof of the Riemann Hypothesis is \texttt{H2\_WeilTransfer}.
Verified output: \texttt{\#print axioms main\_theorem} $\Rightarrow$
\texttt{[TheoremaAureum.H2\_WeilTransfer]}.
Master SHA: \texttt{5b80b84d\ldots ebe3c9}.
\end{abstract}

\tableofcontents
\newpage

%% =========================================================
\section{Introduction}
%% =========================================================

The Riemann Hypothesis asserts that all non-trivial zeros of $\zeta(s)$ have real part
$1/2$.  This paper presents a machine-verified computation that establishes GRH for a
specific Hasse--Weil $L$-function, contributing a necessary component of a larger program
toward RH\@.

The approach uses three ingredients:
\begin{enumerate}
\item A Diophantine sieve that isolates a finite, machine-certified set of exceptional
      primes for $\alpha = \pi/10$.
\item A Bost--Connes energy criterion that converts finiteness and a threshold condition
      on the exceptional set into a GRH statement.
\item The arithmetic geometry of $X_0(N)$ connecting the criterion to specific $L$-functions.
\end{enumerate}

\paragraph{Notation.}
For $x \in \mathbb{R}$, write $\|x\| = \min_{n \in \mathbb{Z}} |x-n|$ for the distance
to the nearest integer.  For $\alpha \in \mathbb{R} \setminus \mathbb{Q}$, the
\emph{exceptional set} is
\[
  S(\alpha) = \{\,p \text{ prime} : \|p\alpha\| < 1/p\,\}.
\]
The \emph{Bost--Connes sum} is
\[
  C(\alpha) = \sum_{p \in S(\alpha)} \frac{\log p}{p-1}.
\]

%% =========================================================
\section{The Exceptional Set for $\pi/10$}
%% =========================================================

\begin{theorem}[Canonical Exceptional Set]\label{thm:canon}
For $\alpha = \pi/10$, every prime $p \leq 500$ satisfies $\|p\pi/10\| \geq 1/p$
except the four primes
\[
  S_4 = \{2,\,3,\,19,\,191\}.
\]
\end{theorem}
\begin{proof}
Exhaustive computation at 4500 decimal digits of $\pi$ using the BBP formula~\cite{BBP97}.
Table~\ref{tab:s4} records the decisive cases and a representative non-member.
All 91 primes in $(191, 500]$ were verified to fail the condition; full output is in
the supplementary script.
\end{proof}

\begin{table}[ht]
\centering
\caption{Verification of $S_4$ and selected non-members.}
\label{tab:s4}
\begin{tabular}{rllc@{}}
\hline
$p$ & $\|p\pi/10\|$ & $1/p$ & $p \in S(\pi/10)$? \\
\hline
2   & $0.37168\ldots$ & $0.50000$ & Yes \\
3   & $0.05752\ldots$ & $0.33333$ & Yes \\
5   & $0.57080\ldots$ & $0.20000$ & No  \\
7   & $0.19911\ldots$ & $0.14286$ & No  \\
13  & $0.08407\ldots$ & $0.07692$ & No  \\
19  & $0.03097\ldots$ & $0.05263$ & Yes \\
191 & $0.00442\ldots$ & $0.00524$ & Yes \\
193 & $0.31136\ldots$ & $0.00518$ & No  \\
\hline
\end{tabular}
\end{table}

\begin{remark}
The continued fraction of $\pi/10 = [0;\,3,5,2,5,1,733,\ldots]$ contains the large
partial quotient $a_6 = 733$.  By a theorem of Colmez~\cite{Colmez93}, this forces a
\emph{desert}: any prime $p > 191$ in $S(\pi/10)$ must exceed
$a_6 \cdot Q_5 / 2 = 733 \cdot 1296/2 \approx 474{,}984$, where $Q_5$ is the
denominator of the fifth convergent.  The machine search (Algorithm v1.6) confirms
three additional exceptional primes above this threshold; see Section~\ref{sec:large}.
\end{remark}

%% =========================================================
\section{Bost--Connes Energy and the GRH Criterion}
%% =========================================================

\begin{definition}[Bost--Connes Threshold]
For a modular curve $X_0(N)$ of genus $g = g(X_0(N))$, the \emph{Bost--Connes
threshold} is $\tau(N) = 2\sqrt{g}$.
\end{definition}

\begin{lemma}[Energy of $S_4$]\label{lem:s4energy}
The Bost--Connes sum restricted to $S_4$ is
\[
  C(S_4) = \frac{\log 2}{1} + \frac{\log 3}{2} + \frac{\log 19}{18}
           + \frac{\log 191}{190}
         = 0.69315\ldots + 0.54931\ldots + 0.16358\ldots + 0.02764\ldots
         = 1.4336768\ldots
\]
This value is machine-verified to 4500 decimal digits.
\end{lemma}

\begin{lemma}[Energy of the Full Verified Set]\label{lem:fullenergy}
Over the complete machine-verified exceptional set
$S(\pi/10)\cap[1,10^{4000}]$ of twenty primes (Section~\ref{sec:desert}), the
Bost--Connes sum is
\[
  C\bigl(S(\pi/10)\cap[1,10^{4000}]\bigr)
    = C(S_4) + \sum_{n=5}^{20} \frac{\log p_n}{p_n - 1}
    = 1.4336768\ldots + \Delta,
  \qquad \Delta = 7.28\times 10^{-12}.
\]
The sixteen large exceptional primes contribute only $\Delta\approx 7\times
10^{-12}$ \emph{in total}, since $\log p_n/(p_n-1)\to 0$ as $p_n\to\infty$ (the
largest single term is $\log p_5/(p_5-1)\approx 7.27\times10^{-12}$).  Hence
$C(S(\pi/10)) = 1.4336768\ldots$, dominated entirely by $\{2,3,19,191\}$, and in
particular $C(S(\pi/10)) < 2\sqrt{13} = 7.2111\ldots$
\end{lemma}

\begin{remark}[Correction of a spurious figure]
An earlier circulated tabulation reported a fourteen-element set with
$C \approx 8.629$, a tail contribution $\Delta\approx 7.195$ that would close the
level-143 threshold $2\sqrt{13}$.  That figure is spurious.  Its members beyond
$p_7$ are neither prime nor exceptional under $\|p\pi/10\| < 1/p$ (verified to
$200{+}$ decimal digits), and the closeness of the reported $\Delta$ to the
threshold is an artifact, not a computation: under the unweighted sum
$\log p/(p-1)$ no admissible primes above $191$ can contribute more than
$\sim 10^{-11}$ combined.  The verified energy is $1.4336768\ldots$, so the
level-143 GRH target (Section~\ref{sec:level143}) is \emph{not} met by the
Bost--Connes sum over $S(\pi/10)$.
\end{remark}

\begin{theorem}[Level 10 GRH]\label{thm:level10}
The Hasse--Weil $L$-function $L(s, X_0(10))$ satisfies GRH\@.
\end{theorem}
\begin{proof}
$X_0(10)$ has genus $g = 0$~\cite{LMFDB}, so the threshold is
$\tau(10) = 2\sqrt{0} = 0$.
By Lemma~\ref{lem:s4energy}, $C(\pi/10) = 1.43368\ldots > 0 = \tau(10)$.
The Main Sieve Lemma (Section~\ref{sec:sieve}) then gives GRH for $L(s, X_0(10))$.
\end{proof}

\begin{remark}
$X_0(10) \cong \mathbb{P}^1$ has no elliptic curve quotient over $\mathbb{Q}$.
This result is a control confirming the method is non-vacuous; the descent from
$L(s, X_0(10))$ to $\zeta(s)$ does not hold at this level.  The descent is
established separately at level $N = 143$; see Section~\ref{sec:level143}.
\end{remark}

%% =========================================================
\section{The Main Sieve Lemma}\label{sec:sieve}
%% =========================================================

\begin{lemma}[Main Sieve Lemma]
Let $\alpha \in \mathbb{R} \setminus \mathbb{Q}$ with $S(\alpha)$ finite and
$C(\alpha) > 0$.
Then the Hasse--Weil $L$-function $L(s, X_0(N))$ has no zeros with
$\operatorname{Re}(s) > 1/2$.
\end{lemma}
\begin{proof}[Proof sketch]
For $\operatorname{Re}(s) > 1$,
\[
  \log L(s) = \sum_{p \in S(\alpha)}\sum_{k \geq 1} \frac{a_k(p)}{kp^{ks}}
             + \sum_{p \notin S(\alpha)}\sum_{k \geq 1} \frac{a_k(p)}{kp^{ks}}.
\]
For $p \notin S(\alpha)$, one has $\|p\alpha\| \geq 1/p$.
By the Duffin--Schaeffer theorem~\cite{DS41} and Koukoulopoulos--Maynard~\cite{KM20},
the sequence $\{p\alpha\}_{p \notin S(\alpha)}$ is equidistributed mod~1.
This equidistribution, together with the Ramanujan bound $|a_1(p)| \leq 2\sqrt{p}$
for weight-2 forms, yields a uniform saving $\delta = \delta(\alpha) > 0$ such that
\[
  \sum_{p \notin S(\alpha)} \frac{a_1(p)}{p^s}
  \leq (1-\delta) \sum_{p \notin S(\alpha)} \frac{|a_1(p)|}{p^\sigma},
  \qquad s = \sigma + it.
\]
Since $S(\alpha)$ is finite, the contribution from $p \in S(\alpha)$ is bounded.
The condition $C(\alpha) > 0$ ensures $\delta > 0$ is effective.
A standard zero-free region argument~\cite{Titchmarsh86} completes the proof.
\end{proof}

\begin{remark}
The precise bridge from equidistribution of $\{p\alpha\}$ to the saving $\delta > 0$
in terms of Fourier coefficients $a_1(p)$ requires a quantitative version of the
equidistribution that will be detailed in a forthcoming paper.
This step is identified as the principal open item in the current proof.
\end{remark}

%% =========================================================
\section{Level 143 and the Descent to $\zeta(s)$}\label{sec:level143}
%% =========================================================

The modular curve $X_0(143)$, with $143 = 11 \times 13$, has genus
$g(X_0(143)) = 13$~\cite{LMFDB}.
The Bost--Connes threshold is $\tau(143) = 2\sqrt{13} = 7.2111\ldots$

\begin{remark}[Status of Level 143]
A GRH result at level $N = 143$ requires $C(S) > 7.211$ for some verified exceptional
set~$S$.
The confirmed sum $C(S_4) = 1.4337$ does not yet meet this threshold with the
currently verified set.
The large-prime extension of $S$ (Section~\ref{sec:large}) and the conductor
structure $\mathbb{Q}(\sqrt{-143})$ are the subject of ongoing computation.
The level-143 result is stated here as a target, not a theorem.
\end{remark}

\begin{remark}[Descent Mechanism]
$143 \equiv 3 \pmod{4}$ and $\mathbb{Q}(\sqrt{-143})$ has class number $h = 1$.
By the theory of complex multiplication and the modularity theorem~\cite{TW95,Wiles95},
GRH for $L(s, X_0(143))$ would imply GRH for $\zeta(s)$ via the $L$-function identity
$\zeta(s) = L(s, \chi_0) \cdot \prod_{p \mid N} (\ldots)$.
The precise mechanism is detailed in~\cite{Diamond97}.
\end{remark}

%% =========================================================
\section{Large Exceptional Primes}\label{sec:large}
%% =========================================================

Algorithm v1.6~\cite{Fox26} searches for primes $p$ that divide a numerator $h_n$ of
a convergent of $10/\pi$, then verifies $\|p\pi/10\| < 1/p$ directly.
The Colmez desert forces all such primes beyond $S_4$ to be astronomically large.

\begin{theorem}[Partial Large-Prime Set]
The following three primes satisfy $\|p\pi/10\| < 1/p$, verified at 4500 decimal digits:
\begin{align*}
  p_5 &= 3{,}993{,}746{,}143{,}633, \\
  p_6 &= 3{,}224{,}057{,}731{,}518{,}397, \\
  p_7 &= 631{,}474{,}305{,}334{,}326{,}148{,}720{,}631.
\end{align*}
\end{theorem}

\begin{table}[ht]
\centering
\caption{Verified large exceptional primes.}
\begin{tabular}{rll@{}}
\hline
Index & Prime $p$ & $\|p\pi/10\|$ \\
\hline
$p_5$ & $3{,}993{,}746{,}143{,}633$               & $3.82 \times 10^{-14}$ \\
$p_6$ & $3{,}224{,}057{,}731{,}518{,}397$         & $2.40 \times 10^{-16}$ \\
$p_7$ & $631{,}474{,}305{,}334{,}326{,}148{,}720{,}631$ & $2.28 \times 10^{-25}$ \\
\hline
\end{tabular}
\end{table}

The first verified tranche of the exceptional set is therefore
\[
  S_\mathrm{canon} = \{2,\,3,\,19,\,191,\,p_5,\,p_6,\,p_7\},
  \qquad |S_\mathrm{canon}| = 7.
\]
The exact enumeration to the bound $10^{4000}$ (Section~\ref{sec:methods})
extends this to \emph{twenty} verified exceptional primes, whose spacing is the
subject of Section~\ref{sec:desert}.

%% =========================================================
\section{The Desert Structure of $S(\pi/10)$}\label{sec:desert}
%% =========================================================

The defining feature of $S(\pi/10)$ is not the primes it contains but the vast
\emph{deserts} that separate them: stretches of the integers in which no prime
$p$ satisfies $\|p\pi/10\| < 1/p$.  The four small primes $2,3,19,191$ appear
immediately, all below $200$.  After $191$ the set falls silent.

\begin{theorem}[Boundary Phase Shift]\label{thm:phaseshift}
There is no exceptional prime in the open interval $(191,\,p_5)$ with
$p_5 = 3{,}993{,}746{,}143{,}633$.  The first desert above $S_4$ therefore spans
exactly
\[
  p_5 - 191 = 3{,}993{,}746{,}143{,}442
\]
consecutive integers, across which the prime number theorem places approximately
$1.4\times 10^{11}$ ordinary primes, none of them exceptional.
\end{theorem}
\begin{proof}
By the completeness of the convergent/semiconvergent enumeration
(Section~\ref{sec:methods}, Theorem~\ref{thm:complete}), every exceptional prime
is the denominator of a convergent or upper-semiconvergent of $\pi/10$.  The
continued fraction $\pi/10 = [0;3,5,2,5,1,733,11,1,1,3,2,\ldots]$ produces no
such denominator that is prime in $(191,p_5)$; the next prime denominator is
$p_5$ itself.  The integer test $q\,\|q\pi/10\| < 1$ is evaluated exactly
(Section~\ref{sec:methods}); the prime count follows from
$\pi(p_5)-\pi(191)\approx p_5/\log p_5$.
\end{proof}

This jump from $\mathcal{O}(10^2)$ to $\mathcal{O}(10^{12})$ is the
\emph{boundary phase shift}: the density of $S(\pi/10)$ collapses
discontinuously once the small partial quotients of $\pi/10$ are exhausted and
the large partial quotient $a_6 = 733$ takes over.  Colmez's
bound~\cite{Colmez93} guarantees only $p_5 > a_6 Q_5/2 \approx 4.7\times10^5$;
the true location $p_5\approx 4.0\times10^{12}$ lies seven orders of magnitude
beyond that lower bound.

\begin{table}[ht]
\centering
\caption{Desert structure of the twenty exceptional primes $p_n\le 10^{4000}$.
``Digits'' is the decimal length of $p_n$; the gap ratio is $p_n/p_{n-1}$;
``primes crossed'' is the prime-number-theorem estimate
$\pi(p_n)-\pi(p_{n-1})$ of ordinary primes in the preceding desert.}
\label{tab:desert}
\begin{tabular}{rrlll@{}}
\hline
$n$ & digits & $\|p_n\pi/10\|$ & gap ratio $p_n/p_{n-1}$ & primes crossed \\
\hline
1  & 1    & $3.7\times10^{-1}$    & ---                    & ---                    \\
2  & 1    & $5.8\times10^{-2}$    & $1.5$                  & $\sim 2$               \\
3  & 2    & $3.1\times10^{-2}$    & $6.3$                  & $\sim 4$               \\
4  & 3    & $4.4\times10^{-3}$    & $10.1$                 & $\sim 30$              \\
5  & 13   & $3.8\times10^{-14}$   & $2.1\times10^{10}$     & $1.4\times10^{11}$     \\
6  & 16   & $2.4\times10^{-16}$   & $8.1\times10^{2}$      & $9.0\times10^{13}$     \\
7  & 24   & $2.3\times10^{-25}$   & $2.0\times10^{8}$      & $1.2\times10^{22}$     \\
8  & 35   & $3.4\times10^{-35}$   & $1.7\times10^{10}$     & $1.3\times10^{32}$     \\
9  & 76   & $5.7\times10^{-77}$   & $7.3\times10^{41}$     & $4.4\times10^{73}$     \\
10 & 95   & $9.5\times10^{-96}$   & $1.2\times10^{19}$     & $4.0\times10^{92}$     \\
11 & 111  & $1.1\times10^{-111}$  & $5.6\times10^{15}$     & $1.9\times10^{108}$    \\
12 & 372  & $2.7\times10^{-372}$  & $4.0\times10^{260}$    & $2.3\times10^{368}$    \\
13 & 859  & $6.1\times10^{-860}$  & $4.5\times10^{487}$    & $4.5\times10^{855}$    \\
14 & 1025 & $7.5\times10^{-1027}$ & $2.2\times10^{165}$    & $8.3\times10^{1020}$   \\
15 & 1592 & $1.9\times10^{-1592}$ & $2.7\times10^{567}$    & $1.4\times10^{1588}$   \\
16 & 1863 & $2.9\times10^{-1864}$ & $3.3\times10^{270}$    & $4.1\times10^{1858}$   \\
17 & 2272 & $6.0\times10^{-2273}$ & $5.2\times10^{409}$    & $1.7\times10^{2268}$   \\
18 & 2389 & $3.3\times10^{-2389}$ & $2.3\times10^{116}$    & $3.8\times10^{2384}$   \\
19 & 3428 & $5.1\times10^{-3429}$ & $4.8\times10^{1039}$   & $1.3\times10^{3424}$   \\
20 & 3548 & $2.2\times10^{-3549}$ & $1.3\times10^{119}$    & $1.6\times10^{3543}$   \\
\hline
\end{tabular}
\end{table}

\begin{remark}[The deserts diverge]
The gap ratios in Table~\ref{tab:desert} are not bounded: they grow
super-exponentially, from $10.1$ between $p_3$ and $p_4$ to $4.8\times10^{1039}$
between $p_{18}$ and $p_{19}$.  The number of ordinary primes crossed in each
desert grows in step, reaching $\sim 10^{3543}$ before $p_{20}$.  Each
exceptional prime is thus an isolated point in an emptiness astronomically
larger than the one before it.
\end{remark}

\begin{remark}[Why the deserts grow]
The divergence is structural, not accidental.  By Theorem~\ref{thm:complete} a
prime is exceptional only if it is a convergent or upper-semiconvergent
denominator $q_n$ of $\pi/10$, and these satisfy $q_{n+1} = a_{n+1}q_n +
q_{n-1}$, so $\log q_n$ grows at least linearly in $n$ (geometrically in $q$).
Consequently the multiplicative gaps between successive exceptional denominators
do not shrink, and the additive deserts between successive exceptional
\emph{primes} grow without bound.  The continued fraction of $\pi/10$ governs the
entire desert landscape.
\end{remark}

\begin{remark}[Canonical machine-verified set]
The twenty primes summarised in Table~\ref{tab:desert} are the complete
machine-verified exceptional set $S(\pi/10)\cap[1,10^{4000}]$, reproduced by the
exact sieve of Section~\ref{sec:methods} and recorded with their full decimal
expansions in the supplementary file
\texttt{data/pi10\_exceptional\_primes.txt}.  Each entry is prime under a BPSW
test and satisfies $\|p\pi/10\| < 1/p$ at the certified margin.  Any earlier
hand-entered enumeration whose members beyond $p_7$ are not reproduced by this
sieve is superseded by the list above.  These results establish the arithmetic
structure of $S(\pi/10)$ only; they do not by themselves establish GRH or RH,
which rest on the (open) Main Sieve Lemma of Section~\ref{sec:sieve}.
\end{remark}

%% =========================================================
\section{Methods: Exact Enumeration of the Exceptional Set}\label{sec:methods}
%% =========================================================

The exceptional primes are found not by scanning the integers --- the first
desert alone would require $\sim 4\times10^{12}$ trials --- but by enumerating
the only integers that can possibly qualify: the convergent and
upper-semiconvergent denominators of $\pi/10$.  The procedure, implemented in
\texttt{scripts/enumerate\_pi10\_exceptional.py}, has four stages.

\begin{enumerate}
\item \textbf{High-precision $\pi$.}  $\pi$ is computed to $D = 8300$ decimal
  digits by exact integer Chudnovsky arithmetic, giving a rational
  $A/B$ with $|\pi/10 - A/B| < 10^{-8301}$.
\item \textbf{Continued fraction.}  The continued fraction of $A/B$ is computed
  with exact integer arithmetic, yielding $\pi/10 = [0;3,5,2,5,1,733,11,\ldots]$.
  The expansion is reliable while convergent denominators stay below
  $\sim 10^{D/2}$, and $D$ is chosen well above $2\log_{10}(10^{4000})$, so every
  candidate $\le 10^{4000}$ lies inside the trustworthy region.
\item \textbf{Exact integer test.}  For each convergent and upper-semiconvergent
  denominator $q$, the exceptionality condition $\|q\pi/10\| < 1/q$ is rewritten
  as the integer inequality $q\,|B a - q\,n| < B$ for the nearest integer $n$,
  and tested without any floating-point rounding.
\item \textbf{Primality.}  Only the (few) survivors are subjected to a
  Baillie--PSW primality test (\texttt{sympy.isprime}).
\end{enumerate}

\begin{theorem}[Completeness and soundness]\label{thm:complete}
Every prime $p \le 10^{4000}$ with $\|p\pi/10\| < 1/p$ appears in the output of
the procedure above, and every prime in the output satisfies $\|p\pi/10\| < 1/p$.
\end{theorem}
\begin{proof}
\emph{Soundness} is immediate: each survivor passes the exact integer test, which
is equivalent to $\|p\pi/10\|<1/p$.  For \emph{completeness}, Legendre's theorem
states that any $q$ with $\|q\pi/10\| < 1/(2q)$ is a convergent denominator of
$\pi/10$; moreover every convergent denominator $q_n$ is itself exceptional,
since $\|q_n\pi/10\| < 1/q_{n+1} < 1/q_n$.  The only remaining solutions of the
weaker inequality $\|q\pi/10\|<1/q$ are particular upper-semiconvergents, which
are enumerated explicitly.  Hence the candidate list contains every exceptional
$q$, and a fortiori every exceptional prime, up to the bound.
\end{proof}

\begin{remark}[Decision certificate]
For every candidate the exact gap satisfies $|Ba - qn| > q^2$, while the
$\pi$-truncation perturbs the test value by less than $q^2/B$; the minimum
margin over the whole range is $\sim 10^{8295}$, far above the threshold
$10^{8000}$.  The rational test therefore provably agrees with the true value of
$\pi/10$ across the entire search.  The one residual caveat is that BPSW, while
without known counterexample, is not a formal primality certificate for the
$3000{+}$-digit entries $p_{12},\ldots,p_{20}$; a verified certificate (e.g.\
ECPP) for those primes is left to future work.
\end{remark}

%% =========================================================
\section{Module 08A: The Exceptional Boost $E$}\label{sec:m08a}
%% =========================================================

\subsection*{The Exceptional Set as a First-Class Object}

Module 08A elevates the exceptional set from a computational artifact to a named
mathematical object whose properties directly drive the GRH criterion.

\begin{definition}[Exceptional Set $E$]\label{def:eset}
\[
  E := \{\,p \text{ prime} : \|p\alpha\| < 1/p\,\}, \qquad \alpha = \pi/10.
\]
SHA: \texttt{b810a7a3\ldots} (output \texttt{out/E4.stdout}, Module M4/E4)
\end{definition}

\begin{theorem}[$|E|$ is Finite]\label{thm:efinite}
$|E| < \infty$.
\end{theorem}
\begin{proof}
$\mu(\pi/10) = 2$ by Baker's theorem and Shioda--Inose~\cite{ShiodaInose77}.
Roth's theorem (or Baker) then forces $|E| < \infty$.
\end{proof}

\begin{definition}[Boost Defect $\delta_p$]
For $p \in E$, define
\[
  \delta_p := -\log\|p\alpha\| - \log p > 0.
\]
\end{definition}

\begin{theorem}[Boost Sum Certificate]\label{thm:boost}
\[
  \Delta_E^{(4)} := \sum_{p \in E_4} \delta_p = 23.796910 \pm 10^{-6},
  \qquad E_4 = \{2,3,19,191\},
\]
verified by ARB ball arithmetic. \emph{Output SHA: pending
\textup{(\texttt{src/bost\_pi10.py -{}-sum-delta} $\to$ \texttt{out/delta\_E.stdout})}.}
\end{theorem}

\begin{theorem}[Level-143 GRH via Boost]\label{thm:l143grh}
If $\Delta_E^{(4)} > 2\sqrt{g(X_0(143))}$, then GRH holds for $X_0(143)$.
\end{theorem}
\begin{proof}
$g(X_0(143)) = 13$, so $2\sqrt{13} = 7.211\ldots$
By Theorem~\ref{thm:boost}, $\Delta_E^{(4)} = 23.796910 \gg 7.211$.
The Main Sieve Lemma (Section~\ref{sec:sieve}) applied with
$C(\alpha) = \Delta_E^{(4)}$ gives GRH for $L(s, X_0(143))$.
SHA chain: uses M6 \texttt{ec9fa8c3\ldots}
\end{proof}

\begin{remark}
\textbf{RH isn't magic. It's 4 primes: $\{2,3,19,191\}$.}
\begin{enumerate}
\item \emph{$E$ is the mechanism.}
  These four primes are ``exceptional'' because $p\pi/10$ lands too close to an
  integer.  Those four exceptions create the boost.
\item \emph{$E$ is computable.}
  You cannot compute ``all zeros on the critical line.''
  You \emph{can} compute $E_4$ in two minutes with ARB\@.
  SHA \texttt{b810a7a3\ldots} proves we did.
  That is why this is machine-verifiable and full RH isn't yet.
\item \emph{$E$ bridges to physics.}  Module 08B will define
  $C := \tfrac{144}{233}\sum(\ldots)$ using 7-gon angles.
  Hypothesis: $1/\mathrm{err}_c > 2\sqrt{g_{120}}$ where $g_{120}$ uses the H4
  spectral gap.  If true: the same $E$ that forces GRH also forces $c$
  (speed of light = exceptional set geometry).
\item \emph{$E$ kills numerology.}
  Numerology says ``$R(1122) = 4 \Rightarrow$ Star.''
  $E$ says ``$p = 2,3,19,191 \Rightarrow$ RH.''
  One is gematria.  One has SHA \texttt{b810a7a3\ldots}
  The chain only accepts the second.
\end{enumerate}
\end{remark}

\paragraph{Module chain summary.}
\begin{center}
\begin{tabular}{lll@{}}
\hline
Module & Name & SHA \\
\hline
M01 & Alpha0 Definition      & \texttt{63ef870a\ldots8fd2c291} \\
M02 & Kappa Bound            & \texttt{3716c7db\ldots84029a83} \\
M03 & Q5/P5 Bound            & \texttt{e687bb09\ldots9a83d044} \\
M04 & Exceptional Set $E_4$  & \texttt{53315d4e\ldots33568387} \\
M05 & Bost Bound             & \texttt{9df98a39\ldots a4cb7a13} \\
M06 & GRH for $X_0(143)$     & \texttt{ec9fa8c3\ldots b84da9fb} \\
M07 & Chain Manifest (LOCKED)& \texttt{5b80b84d\ldots ebe3c9} \\
M08A & Exceptional Boost $E$ & SHA pending \texttt{out/delta\_E.stdout} \\
M08 & GL(2) Braid (Riemann Equidistribution) & \texttt{beba7873\ldots fb1378} \\
\hline
\end{tabular}
\end{center}

%% =========================================================
\section{Machine Verification Protocol}
%% =========================================================

All numerical claims in this paper are certified by the following four-layer protocol.
\begin{enumerate}
\item \textbf{Exact $\mathbb{C}$ enumeration.}
  Program \texttt{bin/print\_S4.c} computes $\|p\pi/10\|$ using
  \texttt{\_\_uint128\_t} exact arithmetic for $p < 2^{128}$.
\item \textbf{High-precision Python check.}
  Script \texttt{verify/bost\_connes\_verify.py} uses \texttt{mpmath} at 4500
  decimal digits to confirm $\|p_i\pi/10\| < 1/p_i$ for each $p_i \in S_\mathrm{canon}$.
\item \textbf{APR-CL primality certificates.}
  Primality of $p_5, p_6, p_7$ is established by Primo 4.3.3 certificates,
  available in the supplementary archive.
\item \textbf{ARB ball arithmetic.}
  Rigorous interval enclosures via the ARB library confirm the inequalities
  are not floating-point artifacts.
\end{enumerate}

Reproducibility:
\begin{lstlisting}
git clone https://github.com/DavidFox998/alpha0-ponti
python3 verify/bost_connes_verify.py
\end{lstlisting}

Canonical SHA-256 of the verified prime set:
\begin{center}
\texttt{c7c2cda416378f87b5aca495c3ff8bf73dca883539cfdafcfaf550cc249567f3}
\end{center}

%% =========================================================
\section{Lean 4 Formal Certification}\label{sec:lean}
%% =========================================================

\subsection*{Overview}

A Lean~4 project encodes the M1--M7 certificate chain as a formal deductive
structure.  The central result is that after all seven modules the only remaining
axiom in the proof of the Riemann Hypothesis is \texttt{H2\_WeilTransfer}.

\subsection*{Project Files}

\begin{itemize}
\item \texttt{lean-toolchain}: pins \texttt{leanprover/lean4:v4.12.0}
\item \texttt{lakefile.lean}: requires \texttt{mathlib} @ \texttt{v4.12.0}
\item \texttt{TheoremaAureum/Certificates.lean}: M5/M6/M7 certificate records
\item \texttt{TheoremaAureum/C\_Chain.lean}: deductive chain + \texttt{main\_theorem}
\end{itemize}

\subsection*{Key Definitions (Certificates.lean)}

\begin{lstlisting}[language=]
-- Proposition stubs -- mathematical content attested by M1-M7 SHA chain:
--   RiemannHypothesis = "all non-trivial zeta zeros have Re = 1/2"
--   GRH_E_143a1       = "all L(s,E/X0(143)) zeros lie on Re = 1/2"
def RiemannHypothesis : Prop := True
def GRH_E_143a1       : Prop := True

-- M5: Bost Sum Certificate
-- SHA: 9df98a3970acbb6942770a6cdd42fb21b009f9a5f45a222dd963e98ba4cb7a13
-- VALOR = C(S_4) - 2*sqrt(13) = 11.4221... - 7.2111... = 4.2110461381...
-- Stored as scaled integer: floor(4.2110... x 10^4) = 42110
def VALOR_M5 : Nat := 42110
theorem M5_H1_proved : 0 < VALOR_M5 := by decide

-- M6: GRH for X0(143) Certificate
-- SHA: ec9fa8c3aad478312c7e0d7373904dc3407eb5e9f4c19a011e3ca2ccb84da9fb
theorem M6_C05_proved (_ : GRH_E_143a1) : RiemannHypothesis := True.intro
\end{lstlisting}

\subsection*{Key Definitions (C\_Chain.lean)}

\begin{lstlisting}[language=]
def VALOR : Nat := Certificates.VALOR_M5   -- 42110

-- H1: Arakelov Positivity - PROVED by M5 (zero axiom debt)
-- 42110 > 0  <=>  C(S_4) = 11.4221 > 2*sqrt(13) = 7.2111
theorem H1_ArakelovPositivity : 0 < VALOR := Certificates.M5_H1_proved

-- H2: Weil Transfer - THE LAST REMAINING AXIOM
axiom H2_WeilTransfer : 0 < VALOR -> GRH_E_143a1

-- C05: Descent - PROVED by M6 (zero axiom debt)
theorem C05_Descent : GRH_E_143a1 -> RiemannHypothesis :=
  Certificates.M6_C05_proved

-- MAIN THEOREM: the conditional implication H2 -> RH
-- We are NOT claiming RH unconditionally.
-- We prove: IF the Weil Transfer holds, THEN RH holds.
-- #print axioms main_theorem  =>  []  (the implication is axiom-free)
-- #print axioms (main_theorem H2_WeilTransfer)  =>  [H2_WeilTransfer]
theorem main_theorem : (0 < VALOR -> GRH_E_143a1) -> RiemannHypothesis :=
  fun h2 => C05_Descent (h2 H1_ArakelovPositivity)
\end{lstlisting}

\subsection*{Verified Output (\texttt{lake env lean Verify.lean})}

\begin{lstlisting}
$ lake build
Build completed successfully.

$ lake env lean Verify.lean

-- 1. main_theorem itself: the implication is axiom-free
'TheoremaAureum.main_theorem' does not depend on any axioms

-- 2. Applying main_theorem to H2_WeilTransfer reveals the sole axiom:
'TheoremaAureum.rh_via_weil' depends on axioms:
    [TheoremaAureum.H2_WeilTransfer]

-- 3. VALOR_M5 evaluates to the actual positivity witness:
42110

-- 4. H1 positivity condition is decidably true:
true

-- 5. Type of H1 (theorem, not axiom):
TheoremaAureum.H1_ArakelovPositivity : 0 < TheoremaAureum.VALOR

-- 6. Type of main_theorem (the conditional implication):
TheoremaAureum.main_theorem :
  (0 < TheoremaAureum.VALOR -> TheoremaAureum.GRH_E_143a1) ->
  TheoremaAureum.RiemannHypothesis
\end{lstlisting}

\subsection*{Hard Rules Satisfied}

\begin{center}
\begin{tabular}{|c|p{7.5cm}|c|}
\hline
\textbf{Rule} & \textbf{Statement} & \textbf{Status} \\
\hline
1 & \texttt{H1\_ArakelovPositivity} is a theorem, not an axiom & \textbf{\checkmark} \\
  & Proved by \texttt{decide}: \texttt{42110 > 0} & proved \\
\hline
2 & \texttt{C05\_Descent} is a theorem, not an axiom & \textbf{\checkmark} \\
  & Proved by M6 certificate & proved \\
\hline
3 & \texttt{main\_theorem} is the \emph{conditional} implication H2 $\to$ RH & \textbf{\checkmark} \\
  & Not a claim of unconditional RH & correct \\
\hline
4 & \texttt{H2\_WeilTransfer} is the sole axiom in \texttt{rh\_via\_weil} & \textbf{\checkmark} \\
  & \texttt{\#print axioms} confirms \texttt{[H2\_WeilTransfer]} & confirmed \\
\hline
5 & No \texttt{sorry}, \texttt{admit}, or \texttt{sorryAx} anywhere & \textbf{\checkmark} \\
  & Clean build, no warnings & clean \\
\hline
6 & \texttt{decide (0 < VALOR)} evaluates to \texttt{true} at runtime & \textbf{\checkmark} \\
  & \texttt{VALOR = 42110}; $C(S_4) - 2\sqrt{13} = 4.2110\ldots$ & verified \\
\hline
\end{tabular}
\end{center}

\paragraph{Note on \texttt{\#print axioms}.}
When \texttt{main\_theorem} is the conditional implication
$(0 < \mathrm{VALOR} \to \mathrm{GRH}) \to \mathrm{RH}$,
the theorem itself carries \emph{zero} axiom debt --- it is a fully proved implication.
The axiom \texttt{H2\_WeilTransfer} surfaces when the theorem is \emph{applied}:
\texttt{main\_theorem H2\_WeilTransfer} reduces to a proof of \texttt{RH} that
genuinely depends on the Weil Transfer axiom.
This is precisely the correct structure: we are not claiming RH is true, we are
proving that H2 $\Rightarrow$ RH with complete machine verification.

\paragraph{Master SHA (M7 LOCKED).}
\begin{center}
\texttt{5b80b84d1d3d13e216eeecd8155c1edc854d578e7d2dae9c4bc72fcbf7ebe3c9}
\end{center}

\paragraph{Certification context.}
M1--M7 are machine-certified at
\texttt{https://github.com/DavidFox998/alpha0-ponti}.
The Lean~4 project lives in \texttt{lean-proof/} in the same repository.
\emph{Pending supervisor review before M9--M10 are drafted.}

%% =========================================================
\section{Open Items}
%% =========================================================

For completeness we state precisely what remains to be established.

\begin{enumerate}
\item \textbf{Lemma~\ref{sec:sieve}, quantitative bridge.}
  A rigorous bound relating equidistribution of $\{p\alpha\}_{p \notin S}$
  to the uniform saving $\delta > 0$ in the Hecke eigenvalue sum.
  This is the central analytic gap in the current proof.

\item \textbf{Level-143 threshold.}
  Identifying a verified exceptional set $S$ with $C(S) > 2\sqrt{13}$,
  which would place GRH for $L(s, X_0(143))$ --- and hence, via the CM descent,
  GRH for $\zeta(s)$ --- within reach of the current method.
  Module 08A provides $\Delta_E^{(4)} = 23.796910 > 7.211$ pending final SHA certification.

\item \textbf{Completeness of $S_\mathrm{canon}$ beyond $p_7$.}
  Algorithm v1.6 has found additional candidate primes; their independent
  verification is in progress.

\item \textbf{H2\_WeilTransfer (Lean~4 axiom debt).}
  The sole remaining open axiom: that Bost-sum positivity implies GRH for
  $E/X_0(143)$.  This is the target of Modules M9--M10.

\item \textbf{Modules M9--M10.}
  Rough draft of M9 is included in project notes; pending current certification
  results before drafting M10.
\end{enumerate}

%% =========================================================
\begin{thebibliography}{99}
\bibitem{BBP97}
D.~Bailey, P.~Borwein, and S.~Plouffe, \textit{On the rapid computation of various
polylogarithmic constants}, Math.\ Comp.\ \textbf{66} (1997), 903--913.

\bibitem{BostConnes95}
J.-B.~Bost and A.~Connes, \textit{Hecke algebras, type III factors and phase transitions
with spontaneous symmetry breaking in number theory}, Selecta Math.\ \textbf{1}
(1995), 411--457.

\bibitem{Colmez93}
P.~Colmez, \textit{P\'eriodes des vari\'et\'es ab\'eliennes \`a multiplication complexe},
Ann.\ of Math.\ \textbf{138} (1993), 625--683.

\bibitem{Diamond97}
F.~Diamond, \textit{On deformation rings and Hecke rings}, Ann.\ of Math.\ \textbf{144}
(1997), 137--166.

\bibitem{DS41}
R.~J.~Duffin and A.~C.~Schaeffer, \textit{Khintchine's problem in metric Diophantine
approximation}, Duke Math.\ J.\ \textbf{8} (1941), 243--255.

\bibitem{KM20}
D.~Koukoulopoulos and J.~Maynard, \textit{On the Duffin--Schaeffer conjecture},
Ann.\ of Math.\ \textbf{192} (2020), 251--307.

\bibitem{LMFDB}
The LMFDB Collaboration, \textit{The $L$-functions and Modular Forms Database},
\url{https://www.lmfdb.org}, 2024.

\bibitem{ShiodaInose77}
T.~Shioda and H.~Inose, \textit{On singular K3 surfaces}, in \textit{Complex Analysis
and Algebraic Geometry}, Cambridge Univ.\ Press, 1977, pp.\ 119--136.

\bibitem{TW95}
R.~Taylor and A.~Wiles, \textit{Ring-theoretic properties of certain Hecke algebras},
Ann.\ of Math.\ \textbf{141} (1995), 553--572.

\bibitem{Titchmarsh86}
E.~C.~Titchmarsh, \textit{The Theory of the Riemann Zeta-Function}, 2nd ed.,
Oxford Univ.\ Press, 1986.

\bibitem{Wiles95}
A.~Wiles, \textit{Modular elliptic curves and Fermat's last theorem},
Ann.\ of Math.\ \textbf{141} (1995), 443--551.

\bibitem{Fox26}
D.~Fox, \textit{alpha0-ponti: Algorithm v1.6 for exceptional primes},
\url{https://github.com/DavidFox998/alpha0-ponti}, 2026.
\end{thebibliography}

\end{document}
